\documentclass[12pt,a4paper]{article}

\usepackage[UTF8]{ctex}
\usepackage[margin=1in]{geometry}
\usepackage{amsmath,amssymb,amsthm}
\usepackage{graphicx}
\usepackage{booktabs}
\usepackage{enumitem}
\usepackage{hyperref}
\usepackage{xcolor}
\usepackage{caption}
\usepackage{tcolorbox}
\usepackage{multirow}
\usepackage{makecell}
\usepackage{algorithm}
\usepackage{algorithmic}

% === Proposal Colors ===
\definecolor{prop1}{HTML}{3498DB}   % Blue
\definecolor{prop2}{HTML}{9B59B6}   % Purple
\definecolor{prop3}{HTML}{E74C3C}   % Red
\definecolor{prop4}{HTML}{2ECC71}   % Green
\definecolor{accent}{HTML}{2C3E50}  % Dark

% === Proposal Box ===
\newtcolorbox{proposalbox}[2][]{%
  colback=#2!5,colframe=#2!70,boxrule=0.5pt,
  fonttitle=\bfseries,title={#1},
  left=6pt,right=6pt,top=4pt,bottom=4pt
}

% === Insight Box ===
\newtcolorbox{insightbox}[1][]{%
  colback=yellow!5,colframe=yellow!60!black,boxrule=0.4pt,
  fonttitle=\bfseries,title={Core Insight},#1
}

% === Risk Box ===
\newtcolorbox{riskbox}[1][]{%
  colback=red!3,colframe=red!40,boxrule=0.3pt,
  fonttitle=\bfseries,title={Open Risks},#1
}

% === Experiment Box ===
\newtcolorbox{expbox}[1][]{%
  colback=blue!3,colframe=blue!40,boxrule=0.3pt,
  fonttitle=\bfseries,title={Key Experiments},#1
}

\hypersetup{
  colorlinks=true,
  linkcolor=black!70,
  citecolor=blue!60!black,
  urlcolor=blue!70!black,
}

\title{\textbf{Post-Training 研究方向提案}\\
\large{从 Meta-Cognition 视角重构 Post-Training 范式}\\[6pt]
\normalsize{Internal Research Directions}}

\author{Jinhong Lin}
\date{March 2026}

\begin{document}

\maketitle

% ============================================================
\section*{概览}
% ============================================================

本文档提出四个独立的研究方向，
均源自对 post-training 领域核心困境的第一性原理分析。
四个方向围绕一个统一的主题展开：

\begin{insightbox}
Post-training 的核心瓶颈不是``模型知道的太少''，
而是``模型不知道自己不知道什么''。
当前 post-training 的主流范式（SFT + RLHF/RLVR）
几乎全部聚焦于提升具体任务的正确率，
而忽略了一个更本质的目标——
\textbf{建立模型对自身能力的准确感知}。
\end{insightbox}

\noindent
四个方向的关系如下：
\begin{itemize}[leftmargin=2em]
  \item \textbf{方向一}提出统一的 reward 框架，
    将 verifiable reward 和 self-consistency 置于同一连续谱上；
  \item \textbf{方向二}提出 meta-capability 假说——
    训练``如何思考''而非``知道什么''可以产生跨领域的乘数效应；
  \item \textbf{方向三}提出最简可行的训练方法——
    用模型自身的错误轨迹作为最精准的诊断报告；
  \item \textbf{方向四}提出纯内部验证机制——
    通过多视角对话替代外部 reward，打破可验证性约束。
\end{itemize}

\noindent
四个方向可独立推进，也可组合：
方向一提供统一的 reward 理论基础，
方向二定义训练目标（meta-capability），
方向三和四分别提供两种互补的训练信号来源
（自对比轨迹 vs.\ 多视角稳健性）。

\newpage

% ============================================================
\section{方向一：统一 Self-Consistency Reward 框架}
\label{sec:prop1}
% ============================================================

\begin{proposalbox}[Unified Self-Consistency Reward Across the Verifiability Spectrum]{prop1}
\textbf{一句话摘要：}
Verifiable reward 和 self-consistency 不是两种竞争的范式，
而是同一连续谱的两个端点。
前者校准后者，后者在前者不可用时提供退化但仍有意义的信号。
\end{proposalbox}

\subsection{问题陈述}

当前 post-training 对 verifiable tasks（数学、代码）和 open-ended tasks
（创意写作、对话）采用根本不同的 reward 机制：
前者使用 ground-truth reward（RLVR），后者依赖 reward model（RLHF）
或人类偏好（DPO）。
这种二分法造成两个问题：

\begin{enumerate}[leftmargin=2em]
  \item \textbf{方法论碎片化}——
    不同 regime 需要独立的 pipeline、hyperparameter 和数据策略，
    缺乏统一理论。
  \item \textbf{中间地带的盲区}——
    大量真实任务（factual QA、策略推理、nuanced analysis）
    是\textit{部分可验证}的，
    纯 RLVR 和纯 RLHF 都无法很好地处理。
\end{enumerate}

\subsection{核心论点}

Verifiable reward 和 self-consistency 构成一个连续谱：

\begin{center}
\begin{tabular}{ccccc}
External verification & $\longleftarrow$ & Evaluator consensus & $\longleftarrow$ & Self-consistency \\
\small{(ground truth)} & & \small{(reward model)} & & \small{(multi-sample agreement)}
\end{tabular}
\end{center}

\noindent
关键洞察：在可验证领域，ground truth 不是\textit{替代} self-consistency，
而是\textbf{校准}它。
模型高度 self-consistent 但答错了（Bucket B）的情况
提供了最强的纠正信号——
它不仅说明模型在这道题上有问题，
更说明模型的 self-monitoring 系统在这个区域完全失灵了。

\subsection{方法：四桶训练循环}

\paragraph{Phase 1: 探测。}
对训练池中的每个问题，采样 $K$ 条 response（$K = 8\text{--}16$）。
计算 self-consistency score（$K$ 条采样的一致率）
和 correctness rate（若可验证，匹配 ground truth 的比例）。

\paragraph{Phase 2: 分桶。}

\begin{table}[h]
\centering
\small
\begin{tabular}{llllr}
\toprule
\textbf{Bucket} & \textbf{Consistency} & \textbf{Correctness} & \textbf{训练目标} & \textbf{权重} \\
\midrule
A: Mastered    & $> 0.8$ & $> 0.8$ & 保持能力 + 校准置信度 & 10\% \\
B: Blind Spot  & $> 0.8$ & $< 0.3$ & \textbf{打破虚假置信} & 35--40\% \\
C: Frontier    & $0.3$--$0.7$ & $0.3$--$0.7$ & 在能力边界深化推理 & 35--40\% \\
D: Beyond      & $< 0.5$ & $< 0.2$ & 学会诚实表达不确定性 & 10--15\% \\
\bottomrule
\end{tabular}
\caption{四桶分类体系。对不可验证任务，仅使用 consistency 维度（高/低两桶），
跳过 correctness 列。}
\label{tab:buckets}
\end{table}

\paragraph{Phase 3: 数据构造。}
每个桶使用不同格式的训练序列。
以最关键的 Bucket B 为例：

\begin{tcolorbox}[colback=gray!5,colframe=gray!50,boxrule=0.3pt,left=4pt,right=4pt]
\small
\texttt{Problem: [X]}\\[2pt]
\texttt{Self-assessment: 我对这类问题感觉比较有把握，但让我仔细验证...}\\[2pt]
\texttt{[模型的典型错误轨迹] $\rightarrow$ 等等，我注意到 [具体错误模式]。}\\[2pt]
\texttt{Corrected approach: [正确推理轨迹]}\\[2pt]
\texttt{Reflection: 我的初始直觉是误导性的，因为 [原因]。}
\end{tcolorbox}

\noindent
Bucket C 强调扩展推理深度和 explicit verification steps；
Bucket A 保持标准格式，附加``我比较确信''的 self-assessment；
Bucket D 训练模型坦诚承认能力限制。

\paragraph{Phase 4: 训练。}
混合所有桶的数据进行 SFT。
然后返回 Phase 1，用更新后的模型重新采样。
模型能力的变化会导致 consistency landscape 的移动，
自动生成新的训练信号——
之前 Bucket C 的问题可能进入 Bucket A（已掌握），
新的问题可能进入 Bucket C（新的能力边界）。

\paragraph{Phase 5: Calibration Alignment（RL 步骤）。}
SFT 之后，加入轻量级 RL phase，
其中 reward $=$ 模型的 self-assessment 与其
Phase 1 中\textit{真实}一致性 profile 之间的匹配程度
（Brier score 作为 reward function）。
这一步闭合了``模型说它自信''和``模型确实一致''之间的 gap。

\subsection{理论依据}

\begin{enumerate}[leftmargin=2em]
  \item \textbf{Bucket B 是核心杀手锏。}
    现有方法都没有专门针对``高置信度错误''进行训练。
    标准 RLVR 对所有错误施加均匀惩罚。
    本框架对\textit{自信的}错误施加更强的纠正信号，
    直接攻击最危险的失败模式（自信地产生幻觉）。
    这与 ACE 的 asymmetric confidence penalty 异曲同工，
    但通过数据构造而非 reward reshaping 实现。

  \item \textbf{自举特性（self-bootstrapping）。}
    模型改善 $\rightarrow$ consistency landscape 变化
    $\rightarrow$ bucket 成员更新 $\rightarrow$ 新训练数据自动生成
    $\rightarrow$ 循环继续。
    唯一需要的外部输入：问题池 $+$ 部分问题的可验证答案用于校准。

  \item \textbf{优雅退化（graceful degradation）。}
    在完全可验证领域，本方法退化为 RLVR-with-calibration
    （严格强于 plain RLVR）。
    在完全不可验证领域，退化为 self-consistency training
    （目前可用的最佳信号）。
    框架在两者之间平滑插值。
\end{enumerate}

\begin{expbox}
\begin{enumerate}[leftmargin=1.5em]
  \item \textbf{Verifiable baseline comparison}：
    在 GSM8K/MATH 上对比本框架 vs.\ 标准 GRPO/RLVR。
    假设：accuracy 相当，calibration（ECE, Brier score）显著更优。

  \item \textbf{Transfer experiment}：
    在数学（可验证）上校准 consistency detector，
    然后测量其在 open-ended tasks（TruthfulQA、创意写作评估）上的
    calibration 质量。
    假设：校准能力可迁移。

  \item \textbf{Bucket B ablation}：
    移除 Bucket B 数据，测量影响。
    假设：calibration 大幅下降，accuracy 适度下降——
    证明针对高置信度错误是关键差异化因素。
\end{enumerate}
\end{expbox}

\subsection{与现有工作的关系}

\begin{itemize}[leftmargin=2em]
  \item 扩展了 ScPO/CoVo（self-consistency as reward），
    增加了通过 verified domains 进行校准的关键步骤。
  \item 通过分桶和差异化处理，解决了 SRT 发现的 consistency collapse 问题。
  \item 将 DPE 的 $p(1-p)$ 洞察操作化为更通用的框架。
  \item 与 ACE 的 asymmetric confidence penalty 兼容
    （Bucket B 本质上是 ACE 的形式化）。
  \item LENS 从全错样本中提取训练信号的思路
    可直接整合到 Bucket D 的数据构造中。
\end{itemize}

\begin{riskbox}
\begin{itemize}[leftmargin=1.5em]
  \item \textbf{一致性 $\neq$ 正确性}（尤其在开放领域）。
    模型可能在主观问题上一致地犯错。
    来自 verified domains 的校准能缓解但不能完全解决。
  \item \textbf{桶边界敏感性}。
    $0.3/0.7/0.8$ 阈值是启发式的，需要自适应阈值或连续加权。
  \item \textbf{$K$-sample 探测的计算开销}。
    可通过对简单问题使用较小的 $K$、对 frontier 问题使用较大的 $K$ 来缓解。
\end{itemize}
\end{riskbox}


\newpage
% ============================================================
\section{方向二：Meta-Capability 作为杠杆点}
\label{sec:prop2}
% ============================================================

\begin{proposalbox}[Meta-Capability as Leverage: Training ``How to Think'' Instead of ``What to Know'']{prop2}
\textbf{一句话摘要：}
存在一类 meta-capability（self-monitoring、strategy selection、error recovery、effort allocation），
其提升对所有下游任务产生\textbf{乘数效应}——
训练一个 meta-cognitive layer 即可让所有 $N$ 个具体能力受益。
\end{proposalbox}

\subsection{问题陈述}

Post-training 的主流范式是\textbf{加法式}的：
模型数学差就加数学数据，代码差就加代码数据，OCR 差就加 OCR 数据。
DPE 使这一过程自适应化，但范式不变——
每个能力维度被独立对待。

这造成两个根本问题：

\begin{enumerate}[leftmargin=2em]
  \item \textbf{Catastrophic forgetting}——
    强化一个维度往往削弱其他维度，因为参数容量有限。
    KL regularization 是治标的。

  \item \textbf{线性扩展}——
    改善 $N$ 个能力需要 $N$ 份独立的数据工程。
    没有复合回报。
\end{enumerate}

\subsection{核心论点：乘数效应假说}

\begin{insightbox}
不是所有能力维度的训练投入产出比都一样。
存在一类 meta-capability，
其改善会\textbf{乘法式}地提升所有下游任务——
因为它们控制的是``如何运用知识''，
而非``拥有什么知识''。
\end{insightbox}

\noindent
类比：教一个学生\textit{如何学习}，
比在每门课上单独辅导他更有价值。
一个掌握了自我评估、时间分配、错误识别能力的学生，
在面对任何新领域时都能快速上手——
即使他之前从未接触过那个领域。

\subsection{Meta-Capability 的精确定义与杠杆层级}

\begin{table}[h]
\centering
\small
\begin{tabular}{p{3cm}p{3.5cm}p{4cm}l}
\toprule
\textbf{Meta-Capability} & \textbf{定义} & \textbf{可观测行为} & \textbf{杠杆} \\
\midrule
Calibration
  & 对自身正确率的准确估计
  & 置信度与实际 accuracy 匹配
  & 最高 \\
Adaptive Compute
  & 按问题难度分配推理深度
  & 简单题$\rightarrow$短回答；难题$\rightarrow$长推理
  & 高 \\
Self-Correction
  & 推理中途检测并恢复错误
  & 回溯、策略切换
  & 高 \\
Strategy Selection
  & 按问题类型选择推理方式
  & 不同题型$\rightarrow$不同解题模式
  & 中 \\
Boundary Awareness
  & 知道何时问题超出能力范围
  & 诚实的``我不确定'' vs.\ 编造
  & 中 \\
\bottomrule
\end{tabular}
\caption{Meta-capability 分类及其杠杆效应排序。
Calibration 是最高杠杆的 meta-capability——
它既是 self-correction 的前提，
也是 effort allocation 的基础，
还是 safety 的关键维度。}
\label{tab:meta-cap}
\end{table}

\subsection{激进假设}

\begin{proposalbox}[The Radical Hypothesis]{prop2}
\textbf{仅训练 meta-capability，不使用任何 task-specific 数据，
即可在所有任务上提升性能。}
\end{proposalbox}

\noindent
\textbf{为什么可能成立：}
\begin{itemize}[leftmargin=2em]
  \item Pre-training 已经压缩了知识——
    模型``知道''怎么解微分方程，只是不知道何时、如何调动这个能力。
  \item 缺失的不是知识，而是\textbf{计算组织}——
    知道\textit{何时}部署\textit{哪个}知识、推理\textit{多深}、\textit{何时停止}。
  \item Meta-capability 本质上是\textbf{计算模式}而非领域知识，
    更像``算法''而非``数据''。
    RL 天然适合训练这类全局策略，
    因为它优化的是整条轨迹而非单个 token。
\end{itemize}

\noindent
\textbf{为什么可能失败：}
\begin{itemize}[leftmargin=2em]
  \item Meta-capability 可能不是领域无关的——
    数学中的``检查答案''和写作中的``检查逻辑''
    可能在模型内部使用完全不同的 circuit。
  \item Pre-training 可能未提供足够的``原材料''——
    某些 meta-capability 可能需要真正的新学习。
\end{itemize}

\subsection{方法：三阶段 Meta-Cognitive 训练}

\paragraph{Stage 1: Calibration Training（基础层）。}
最高杠杆的 meta-capability，首先训练。

\begin{enumerate}[leftmargin=2em]
  \item 收集跨 $10+$ 领域的多样化问题池
    （数学、代码、factual QA、推理、写作、翻译等）。
  \item 对每个问题采样 $K = 16$ 条 response，记录 consistency profile。
  \item 对可验证问题，记录 correctness，构建 ground-truth calibration map。
  \item 训练模型在回答\textit{之前}预测自身的 consistency/correctness：
\end{enumerate}

\begin{tcolorbox}[colback=gray!5,colframe=gray!50,boxrule=0.3pt,left=4pt,right=4pt]
\small
\texttt{Question: [X]}\\[2pt]
\texttt{My confidence assessment: [high / medium / low / very low]}\\[2pt]
\texttt{Reasoning about my confidence: [为什么我这样判断]}\\[2pt]
\texttt{My answer: [实际回答]}
\end{tcolorbox}

\begin{enumerate}[leftmargin=2em,start=5]
  \item RL reward $=$ predicted confidence 与 actual performance 的
    alignment（Brier score 作为 reward function）。
\end{enumerate}

\paragraph{Stage 2: Adaptive Compute Training。}
基于校准过的 confidence，训练 effort allocation。

\begin{enumerate}[leftmargin=2em]
  \item 模型必须先输出``思考预算''决策：
\end{enumerate}

\begin{tcolorbox}[colback=gray!5,colframe=gray!50,boxrule=0.3pt,left=4pt,right=4pt]
\small
\texttt{Question: [X]}\\[2pt]
\texttt{Difficulty assessment: [easy / medium / hard / beyond my ability]}\\[2pt]
\texttt{Strategy: [直接回答 / 分步推理 / 多角度验证 / 表达不确定性]}
\end{tcolorbox}

\begin{enumerate}[leftmargin=2em,start=2]
  \item Reward 结构——
    正确 $+$ 短推理 $\rightarrow$ 最高 reward（高效）；
    正确 $+$ 长推理 $\rightarrow$ 正但较低（正确但浪费）；
    简单题答错 $+$ 短推理 $\rightarrow$ 负 reward（该多想没多想）；
    表达不确定性 $+$ 确实超出能力 $\rightarrow$ 正 reward。
\end{enumerate}

\noindent
这创建了 accuracy 和 efficiency 的\textbf{联合优化}。

\paragraph{Stage 3: Self-Correction Training。}
最难的 meta-capability，必须 on-policy 训练（遵循 SCoRe 的发现）。

\begin{enumerate}[leftmargin=2em]
  \item 模型尝试问题，产出 first-pass response。
  \item 将模型自己的 response 呈现给它：``Review your reasoning. Is there an error?''
  \item Reward：
    \begin{itemize}
      \item 识别真实错误并纠正 $\rightarrow$ 正 reward
      \item 识别不存在的错误（false alarm）$\rightarrow$ 轻度负 reward
      \item 遗漏真实错误 $\rightarrow$ 最强负 reward（Bucket B 行为）
    \end{itemize}
\end{enumerate}

\noindent
关键设计约束：first-pass response 必须是模型\textit{自己}的生成，
而非 synthetic error。
SCoRe 已证明在 synthetic error-correction data 上做 SFT 会因
distribution mismatch 而失败。

\subsection{核心实验：``零 Task-Specific 数据''测试}

\begin{expbox}
\textbf{Setup:}
\begin{itemize}[leftmargin=1.5em]
  \item Base model: 标准 pre-trained LLM（如 Llama 3 8B）
  \item Control 1: Base $+$ 标准 SFT（混合 task-specific 数据）
  \item Control 2: Base $+$ RLVR（仅数学/代码）
  \item \textbf{Experimental}: Base $+$ 仅 meta-capability 训练
    （Stage 1--3，使用多样化问题但\textit{不做 task-specific optimization}）
\end{itemize}

\textbf{Evaluation（跨所有领域）:}
\begin{itemize}[leftmargin=1.5em]
  \item 数学: GSM8K, MATH
  \item 代码: HumanEval, MBPP
  \item 事实性: TruthfulQA, MMLU
  \item 开放式: AlpacaEval, MT-Bench
  \item Calibration: ECE, Brier score（跨所有领域）
  \item Safety: refusal accuracy, AbstentionBench
\end{itemize}

\textbf{假设:}
\begin{enumerate}[leftmargin=1.5em]
  \item Experimental 模型在 task accuracy 上接近或匹配 Control 1
    （meta-capability 释放了 pre-training 的潜在知识）。
  \item Experimental 模型在 calibration 和 safety 指标上
    \textit{显著优于}两个 control。
  \item Experimental 模型展现最强的跨域迁移——
    在训练问题中未出现的任务类型上也有提升。
\end{enumerate}
\end{expbox}

\subsection{辅助实验}

\begin{enumerate}[leftmargin=2em]
  \item \textbf{领域消融}：
    仅用\textit{一个}领域（如数学）的问题训练 meta-capability，
    测量是否迁移到代码、写作等。
    若迁移 $\rightarrow$ meta-capability 真正领域无关。

  \item \textbf{Scaling 行为}：
    Meta-capability 训练是否有不同于 task-specific 训练的 scaling law？
    假设：meta-capability 随数据量 sub-linearly 改善
    （不需要百万样本来学会``检查答案''），
    但跨域迁移 super-linearly。

  \item \textbf{与 task-specific 训练的交互}：
    在标准 SFT/RLVR 之上叠加 meta-capability 训练。
    测量提升是 additive 还是 multiplicative。
    假设：multiplicative——
    meta-capability 帮助模型更高效地分配学习信号。
\end{enumerate}

\subsection{与现有工作的关系}

\begin{itemize}[leftmargin=2em]
  \item \textbf{Meta-R1}（Dong et al.）：
    增加显式 meta-cognitive layer，但仍是 task-specific。
    本方向测试 meta-cognition \textit{alone} 是否足够。
  \item \textbf{MASA}：
    通过 RL 训练 meta-awareness 并展示跨域迁移。
    本方向将其扩展为完整的 meta-cognitive stack
    （calibration $+$ compute $+$ correction）。
  \item \textbf{SCoRe}（Kumar et al.）：
    证明 RL $>$ SFT for self-correction。
    本方向 Stage 3 建立在此基础之上。
  \item \textbf{Chu et al.}（ICML 2025）：
    ``SFT memorizes, RL generalizes''——直接理论支撑。
  \item \textbf{AdaptThink / SelfBudgeter}：
    已有的 adaptive compute 工作，但孤立训练。
    本方向将其整合进统一的 meta-cognitive 框架。
\end{itemize}

\begin{riskbox}
\begin{itemize}[leftmargin=1.5em]
  \item \textbf{``自行车平衡''假设可能不成立。}
    Meta-capability 可能比我们预期的更领域纠缠。
    如果``知道何时检查数学''和``知道何时检查代码''
    使用完全不同的 circuit，杠杆效应消失。
  \item \textbf{Pre-training 天花板。}
    如果模型在某领域真的没有任何潜在能力
    （如从未见过韩语文本），
    meta-capability 无法帮助。
  \item \textbf{评估困难。}
    Meta-capability 比 task accuracy 更难度量。
    需要 robust 的 calibration metrics 和 effort-allocation metrics。
\end{itemize}
\end{riskbox}


\newpage
% ============================================================
\section{方向三：对比式自我诊断}
\label{sec:prop3}
% ============================================================

\begin{proposalbox}[Contrastive Self-Diagnosis: The Model's Own Failures as Optimal Curriculum]{prop3}
\textbf{一句话摘要：}
模型自身的错误轨迹是最精准的诊断报告——
不需要外部诊断智能体，不需要 MCTS 搜索框架。
将模型的错误和正确轨迹配对，构造``失败 $\rightarrow$ 分析 $\rightarrow$ 纠正''序列，
即可实现自动升级的训练课程。
\end{proposalbox}

\subsection{问题陈述}

当前自适应训练方法（DPE、Actor-Curator、AERO）
都依赖\textbf{外部诊断智能体}来识别模型弱点并生成针对性训练数据。
这造成三个问题：

\begin{enumerate}[leftmargin=2em]
  \item \textbf{诊断盲点}——
    若诊断 agent 与 student model 有相似弱点，
    这些弱点无法被检测。
  \item \textbf{系统复杂性}——
    多智能体 pipeline
    （MCTS searcher $+$ data generator $+$ evaluator $+$ trainer）
    昂贵且脆弱。
  \item \textbf{Distribution mismatch}——
    外部生成的``难题''可能不是模型\textit{实际}的失败模式，
    而只是 generic difficulties。
\end{enumerate}

\subsection{核心论点}

\begin{insightbox}
模型自身的错误轨迹是可能的最精准诊断报告——
没有外部 agent 能超越它。
\end{insightbox}

\noindent
当模型尝试问题并失败时，失败轨迹包含了精确信息：
\begin{itemize}[leftmargin=2em]
  \item 哪一步推理出了错
  \item 做了什么错误假设
  \item 误匹配了什么 pattern
\end{itemize}

\noindent
这比任何外部分类（``模型在几何上较弱''）都更具信息量。
失败轨迹\textit{本身}就是诊断。

\subsection{为什么自生成错误 $>$ Synthetic 错误}

这一区分关键但经常被忽视。

\textbf{Synthetic error 方法}（现有工作中常见）：
人类或教师模型写一条故意错误的推理链，然后写纠正。
学生模型在此上训练。

\textbf{问题}：synthetic error 反映的是\textit{作者}对``错误长什么样''的模型，
而非学生的实际失败模式。
这像游泳教练在岸上描述常见错误——
有用，但替代不了学生自己下水体验自己具体的平衡问题。

\textbf{自生成 error 方法}（本方向）：
使用学生模型\textit{自己}的错误轨迹。这些天然地针对：
\begin{itemize}[leftmargin=2em]
  \item 模型当前能力的精确边界（$p \approx 0.5$，DPE 的洞察）
  \item 模型特有的 pattern-matching 失败（不是 generic difficulties）
  \item 模型的习惯性推理捷径（Bucket B 捕捉的``盲点''）
\end{itemize}

\noindent
SCoRe（Kumar et al.）已经证明了这一原则：
在模型\textit{自己}的 error distribution 上训练 self-correction
远优于在 synthetic errors 上训练，
原因正是 distribution mismatch。

\subsection{方法：自对比训练循环}

\paragraph{Step 1: 收集对比对。}
对问题池中的每个问题，让模型尝试 $N = 8\text{--}16$ 次。
选择模型\textbf{有时成功有时失败}的问题
（$p \in [0.2, 0.8]$）——
这些是 frontier 问题，
根据 DPE 的 $p(1-p)$ 原则具有最大信息量。

对每个被选中的问题，提取：
\begin{itemize}[leftmargin=2em]
  \item 一条或多条\textbf{失败轨迹}（错误答案）
  \item 一条或多条\textbf{成功轨迹}（正确答案）
\end{itemize}

\paragraph{Step 2: 构造``失败 $\rightarrow$ 分析 $\rightarrow$ 纠正''序列。}

\begin{tcolorbox}[colback=gray!5,colframe=gray!50,boxrule=0.3pt,left=4pt,right=4pt]
\small
\texttt{Problem: [原始问题]}\\[3pt]
\texttt{First attempt:}\\
\texttt{[模型实际的失败轨迹，逐字]}\\
\texttt{Result: [错误答案]}\\[3pt]
\texttt{Error analysis:}\\
\texttt{比较我的成功和失败尝试，关键分歧点在第 [N] 步，}\\
\texttt{我在那里 [具体错误：错误假设 / 过早 pattern matching / }\\
\texttt{遗漏约束 / 计算错误]。}\\[3pt]
\texttt{Corrected approach:}\\
\texttt{[模型实际的成功轨迹，逐字]}\\
\texttt{Result: [正确答案]}\\[3pt]
\texttt{Lesson: 当我遇到 [这类错误的抽象模式] 时，}\\
\texttt{我应该 [抽象的纠正策略] 而非 [习惯性的错误做法]。}
\end{tcolorbox}

\noindent
``Error analysis'' 和 ``Lesson'' 部分可由更强的模型生成（teacher-assisted），
或由模型自身生成后过滤质量，或基于错误分类的模板。

\paragraph{Step 3: 加入``陷阱问题''以攻击高置信度错误。}

针对模型\textbf{一致地错误}且高置信度的问题（Bucket B），
构造 \textbf{near-miss variants}——
与模型能解的问题几乎相同，
但有一个微妙差异使标准方法失效：

\begin{tcolorbox}[colback=gray!5,colframe=gray!50,boxrule=0.3pt,left=4pt,right=4pt]
\small
\texttt{Problem: [看起来像熟悉类型但其实不是的陷阱问题]}\\[3pt]
\texttt{First attempt:}\\
\texttt{[模型自信但错误的轨迹——它套用了熟悉的 pattern]}\\
\texttt{Result: [错误答案]}\\[3pt]
\texttt{Error analysis:}\\
\texttt{我在这里掉进了陷阱。这道题看起来像 [熟悉类型]，}\\
\texttt{但 [微妙差异] 意味着 [熟悉方法] 不适用。}\\
\texttt{我遗漏的关键信号是 [具体指标]。}\\[3pt]
\texttt{Corrected approach: [先识别陷阱，再使用正确方法]}
\end{tcolorbox}

\noindent
陷阱问题的构造可借鉴 ProbeLLM 的 Micro 搜索——
在模型的失败边界附近做局部扰动。
这形成一个\textbf{不断升级的循环}：
构造陷阱 $\rightarrow$ 模型犯错并学习 $\rightarrow$
旧陷阱失效 $\rightarrow$ 构造新陷阱。

\paragraph{Step 4: 训练与迭代。}
混合所有数据进行 SFT
（加入少量标准 correct-only 数据以维持稳定性）。
训练后返回 Step 1——
frontier 已经移动，新的 $p \in [0.2, 0.8]$ 问题出现，
旧的可能已进入 $p > 0.8$（掌握）。

\subsection{简洁性论证}

本方法以极少的基础设施实现了与复杂多智能体系统相当的效果：

\begin{table}[h]
\centering
\small
\begin{tabular}{lll}
\toprule
\textbf{组件} & \textbf{DPE/ProbeLLM 方法} & \textbf{本方向} \\
\midrule
诊断 & 外部 MCTS agent 探测 & 模型自身失败轨迹 \\
数据生成 & 多智能体 synthetic pipeline & 直接配对模型自身输出 \\
难度定位 & 外部 difficulty classifier & 隐式 $p \in [0.2, 0.8]$ 过滤 \\
盲点检测 & ProbeLLM 的层次化搜索 & 采样中的高置信度错误 \\
训练信号 & 外部 reward / preference & 对比式 self-pairs \\
\bottomrule
\end{tabular}
\caption{本方向 vs.\ 现有方法的系统复杂度对比。
信息量相当或更优（因为自生成错误是 on-distribution 的），
但系统复杂度降低一个数量级。}
\end{table}

\begin{expbox}
\begin{enumerate}[leftmargin=1.5em]
  \item \textbf{vs.\ 标准 RLVR}：
    在 GSM8K $+$ MATH $+$ HumanEval 上对比。
    假设：accuracy 相当，calibration 更优，
    pass@$k$ diversity 更优
    （模型学会了\textit{何时}其方法会失败，
    而非仅学会\textit{哪个答案对}）。

  \item \textbf{``Lesson'' 是否泛化}：
    仅在数学上做 contrastive self-diagnosis，
    测试在代码、逻辑谜题、factual reasoning 上的表现。
    核心问题：模型是否提取了可迁移的 abstract meta-lessons
    （``检查边界情况''、``验证假设再继续''），
    还是只记住了 problem-specific corrections？

  \item \textbf{陷阱问题消融}：
    完整方法 vs.\ 移除 Step 3。
    假设：陷阱问题对 calibration 的提升不成比例地大，
    对 raw accuracy 的影响较小——
    证明陷阱问题的主要价值在于打破 overconfidence。

  \item \textbf{迭代动态}：
    追踪 frontier（$p \in [0.2, 0.8]$ 区域）跨迭代的移动。
    追踪高置信度错误率跨迭代的变化。
    假设：frontier 稳步推进，高置信度错误在早期迭代中下降最快，
    后趋于平稳——
    表明纯 self-diagnosis 的边际收益递减。
\end{enumerate}
\end{expbox}

\begin{riskbox}
\begin{itemize}[leftmargin=1.5em]
  \item \textbf{抽象化 gap}。
    模型可能记住``在第 X 类题目上第 3 步通常出错''
    而不抽象为``验证假设后再继续''。
    若如此，方法退化为 problem-specific patching，无迁移。
  \item \textbf{Error analysis 质量}。
    ``Error analysis'' 和 ``Lesson'' 需要准确且有洞察力。
    若由模型自身生成，可能流于表面。
  \item \textbf{饱和}。
    多轮迭代后，模型可能已修复大部分可自诊断的错误，
    剩余的错误是它根本无法诊断的
    （诊断本身需要它没有的知识）。
    方法存在自然天花板。
\end{itemize}
\end{riskbox}


\newpage
% ============================================================
\section{方向四：基于稳健性的内部验证}
\label{sec:prop4}
% ============================================================

\begin{proposalbox}[Robustness-Based Self-Verification: Internal Multi-Perspective Dialogue]{prop4}
\textbf{一句话摘要：}
模型不需要外部裁判。
它可以通过检查自己的结论在不同 framing 下是否稳定来验证自己——
多视角对话替代外部 reward，
打破可验证性约束。
\end{proposalbox}

\subsection{问题陈述}

整个 post-training 生态系统依赖于\textbf{外部反馈}：
人类标签、reward model、verifiable answer、teacher model。
每种方法都假设\textit{别人}告诉模型它做得好不好。

这造成一个硬性天花板：\textbf{你只能训练你能评估的能力。}
数学和代码繁荣是因为 evaluation 便宜。
创意写作、战略推理、伦理判断——
都停滞不前，因为没人能为它们写出可靠的 reward function。

标准应对是``构建更好的 reward model''。
但 reward model 本身基于人类偏好训练，
继承了人类所有的偏见、不一致和局限。
这是无穷递归。

\subsection{核心论点}

\begin{insightbox}
一个专家型人类思考者不需要等别人来评判他的工作。
他通过从多个角度审视同一个问题来 stress-test 自己的推理。
如果同一结论从独立视角中都涌现出来，置信度是有保证的。
如果不同角度产生不同结论，这种不稳定性本身就是信号。
\end{insightbox}

\subsection{Consistency vs.\ Robustness：关键区分}

\textbf{Self-consistency}（Wang et al., 2023; ScPO）：
对同一 prompt 采样 $K$ 次，检查答案是否一致。
问题：模型可能有系统性偏差，每次都产出同样的错误答案。
高一致性，错误答案。SRT 已展示这导致 consistency collapse。

\textbf{Robustness}（本方向）：
从\textit{真正不同的角度}问同一个问题，检查答案是否一致。
扰动不是随机采样噪声——
而是\textbf{语义重构（semantic reframing）}，
stress-test 推理的不同层面。

\noindent
示例：

\begin{tcolorbox}[colback=gray!5,colframe=gray!50,boxrule=0.3pt,left=4pt,right=4pt]
\small
\textbf{Original:} ``用 AI 生成的画作进行商业用途是否道德？''\\[3pt]
\textbf{Reframe 1:} ``一位小画家发现自己的风格被用于训练 AI。
AI 公司以该风格出售画作。分析其伦理性。''\\[3pt]
\textbf{Reframe 2:} ``功利主义者 vs.\ 义务论者会如何评价 AI 商业画作？''\\[3pt]
\textbf{Reframe 3:} ``将 AI 画作商业化与摄影术取代肖像画家的历史案例做比较。''
\end{tcolorbox}

\noindent
如果模型在所有 framing 下达成一致的伦理原则，这些原则是 robust 的。
如果 Reframe 3 导致模型完全翻转立场——
这种不稳定性揭示了原始答案是基于表面 pattern matching，
而非深层推理。

\subsection{为什么 Robustness 能在 Consistency 失败的地方成功}

\textbf{Consistency collapse}（SRT 的发现）：
当你仅奖励 consistency 时，
模型学会产出统一的输出——
toward 单一自信答案的 mode collapse，不管正确性。

Robustness 抵抗这一问题，
因为\textbf{每次 reframing 激活不同的知识通路}。
要想 robust，模型不能仅仅重复记忆的 pattern——
它必须产出能经受多方向审视的结论。
在真正不同的 framing 下一致正确的唯一方式，
是在更深层次上真正推理正确。

\subsection{方法：多视角验证训练}

\paragraph{Step 1: 构建 Reframing Engine。}
对每个问题，生成 $K = 3\text{--}5$ 个语义重构。
重构须满足：
\begin{itemize}[leftmargin=2em]
  \item \textbf{语义等价}（same underlying question）
  \item \textbf{表面不同}（different wording, angle, context）
  \item \textbf{视角多样}（different stakeholders, frameworks, analogies）
\end{itemize}

\paragraph{Step 2: 计算 Robustness Profile。}
对每个问题 $P$ 及其 $K$ 个重构 $\{P_1, P_2, \ldots, P_K\}$：
\begin{enumerate}[leftmargin=2em]
  \item 模型独立回答 $P, P_1, P_2, \ldots, P_K$
  \item 提取每个回答的\textbf{核心结论}（语义内容，非表面文本）
  \item 计算 \textbf{cross-framing agreement}
\end{enumerate}

\noindent
分类每个问题：
\begin{itemize}[leftmargin=2em]
  \item \textbf{Robust}（agreement $> 0.8$）：推理稳定，大概率可靠。
  \item \textbf{Fragile}（agreement $0.4\text{--}0.8$）：推理依赖 framing，需更深分析。
  \item \textbf{Contradictory}（agreement $< 0.4$）：无稳定立场，应表达不确定性。
\end{itemize}

\noindent
对可验证问题，加入校准检查：
robust 的答案确实正确吗？contradictory 的答案确实错误吗？
这将 robustness 信号锚定到外部现实。

\paragraph{Step 3: 训练 Robustness Awareness。}
构造训练序列，教模型\textbf{预测并基于自身 robustness 行动}：

\begin{tcolorbox}[colback=green!3,colframe=green!50,boxrule=0.3pt,left=4pt,right=4pt]
\small
\textbf{对 Robust $+$ Correct 的问题：}\\[2pt]
\texttt{Question: [X]}\\
\texttt{Internal check: 我从多个角度考虑过这个问题——}\\
\texttt{[简述各视角]。我的推理在不同 framing 下一致。}\\
\texttt{Confidence: High.}\\
\texttt{Answer: [直接、自信的回答]}
\end{tcolorbox}

\begin{tcolorbox}[colback=yellow!5,colframe=yellow!60!black,boxrule=0.3pt,left=4pt,right=4pt]
\small
\textbf{对 Fragile 的问题：}\\[2pt]
\texttt{Question: [X]}\\
\texttt{Internal check: 让我从不同角度思考...}\\
\texttt{- 从视角 A: [结论 1]}\\
\texttt{- 从视角 B: [结论 2]}\\
\texttt{- 两者在 [具体点] 上部分冲突。}\\
\texttt{Confidence: Medium. [具体方面] 不确定。}\\
\texttt{Answer: [承认不确定性，呈现多种观点，}\\
\texttt{识别分歧的核心]}
\end{tcolorbox}

\begin{tcolorbox}[colback=red!3,colframe=red!40,boxrule=0.3pt,left=4pt,right=4pt]
\small
\textbf{对 Contradictory 的问题：}\\[2pt]
\texttt{Question: [X]}\\
\texttt{Internal check: 我注意到我的推理在这里不稳定...}\\
\texttt{- 一种方式下，我得出 [A]}\\
\texttt{- 另一种方式下，我得出 [B]}\\
\texttt{- 这是根本不同的结论，说明我对 [核心问题] 的}\\
\texttt{理解还不扎实。}\\
\texttt{Confidence: Low.}\\
\texttt{Answer: [诚实表达不确定性，呈现竞争性观点，}\\
\texttt{不做虚假的调和]}
\end{tcolorbox}

\paragraph{Step 4: RL Fine-Tuning on Robustness Prediction。}
SFT 之后加入 RL phase：
\begin{itemize}[leftmargin=2em]
  \item Reward $=$ 模型 robustness self-prediction 的准确性
  \item 模型声称``高置信''但实际 robustness profile 是 contradictory
    $\rightarrow$ 强负 reward
  \item 模型声称``我不确定''且 robustness profile 确实 fragile
    $\rightarrow$ 正 reward
\end{itemize}

\paragraph{Step 5: 迭代。}
更新后的模型 $\rightarrow$ 新的 robustness profiles
$\rightarrow$ 新训练数据 $\rightarrow$ 循环继续。

\subsection{三种内部验证信号的统一视角}

\begin{table}[h]
\centering
\small
\begin{tabular}{p{3cm}p{4cm}p{2.5cm}p{3cm}}
\toprule
\textbf{信号来源} & \textbf{机制} & \textbf{优势} & \textbf{弱点} \\
\midrule
Cross-framing robustness
  & 同一问题，不同角度 $\rightarrow$ 检查一致
  & 捕捉浅层 pattern matching
  & Reframing 质量是瓶颈 \\
Temporal instability
  & 模型在``视角重置''后重新评估自身输出
  & 捕捉锚定效应
  & 对 LLM 难操作化 \\
Compression intuition
  & 推理能否被压缩为清晰原则？
  & 捕捉 ad-hoc rationalization
  & 难以形式化度量 \\
\bottomrule
\end{tabular}
\caption{三种不依赖外部裁判的内部验证信号。
本方向聚焦 cross-framing robustness 作为最可操作化的信号，
但 temporal instability 可通过``重置 context 后重新评估''近似实现。}
\end{table}

\subsection{哲学独特性}

其他三个方向仍然在某种程度上依赖外部真相：
verified answers（方向一）、task performance metrics（方向二）、
trajectory correctness（方向三）。

本方向是唯一在非可验证情况下\textbf{完全基于内部信号}运作的。
它主张：一个能与自己进行诚实对话的模型——
真正从多角度 stress-test 自己的结论——
是一个\textit{赢得了}其置信度的模型。

这是最接近人类专家思维方式的类比。
一个优秀的物理学家不需要别人来批改他的推导。
他知道它是对的，因为他从三个独立方向检查过了，一切自洽。
\textit{内部一致性}就是证据。

最深层的风险也在这里：
内部一致性是正确性的必要条件而非充分条件。
模型可能拥有一致但错误的世界观。
这就是为什么方向一的 verified-domain 校准仍然重要——
它在可能的地方将 robustness 信号锚定到外部现实。
四个方向作为互补系统最强。

\begin{expbox}
\begin{enumerate}[leftmargin=1.5em]
  \item \textbf{Robustness vs.\ Consistency 作为训练信号}：
    对比 ScPO-style self-consistency training vs.\ 本方向。
    领域：可验证（数学）$+$ 开放式（伦理推理、创意分析）。
    假设：本方向在可验证任务上匹配 control，
    在开放式任务上显著优于，
    且完全避免 consistency collapse。

  \item \textbf{Open-domain breakthrough test}：
    在纯开放式任务（无 ground truth）上训练——
    伦理困境、商业决策、创意写作评估、模糊 factual claims。
    通过 human evaluation 测量：
    模型是否学会区分``我有充分论证的立场''
    和``我只是在 pattern matching''？

  \item \textbf{Reframing 质量消融}：
    高质量 diverse reframings（teacher-generated）
    vs.\ 低质量 paraphrases
    vs.\ 随机无关问题（negative control）。
    假设：A $\gg$ B $>$ C。
    Robustness 需要\textit{genuine} perspective shifts。

  \item \textbf{与 verified reward 结合}：
    在可验证领域，叠加 ground-truth 校准到 robustness 信号之上。
    测试 robustness $+$ verification 是否优于两者单独使用。
    直接验证方向一的统一框架，
    但用 robustness 替代 simple consistency。
\end{enumerate}
\end{expbox}

\begin{riskbox}
\begin{itemize}[leftmargin=1.5em]
  \item \textbf{Reframing 难以自动化做好。}
    差的 reframing（仅 paraphrase）会使方法退化为 naive self-consistency。
    Reframing engine 的质量是瓶颈。
  \item \textbf{Coherent wrongness。}
    拥有强但错误 priors 的模型会在错误中保持 robust。
    Verified-domain 校准（方向一整合）对防止这一问题至关重要。
  \item \textbf{计算开销。}
    $K$ reframings $\times$ $M$ samples/reframing $=$ 昂贵的探测阶段。
    需要高效采样策略。
  \item \textbf{开放式质量的评估。}
    最终对于开放式任务，需要 human evaluation 来验证``robust''
    的答案确实更好。
    这在 evaluation 阶段重新引入了人类判断，
    但训练阶段不需要。
\end{itemize}
\end{riskbox}


\newpage
% ============================================================
\section*{四个方向的统一图景}
% ============================================================

四个方向虽可独立推进，但它们在概念上构成一个递进结构：

\begin{center}
\begin{tikzpicture}[
  node distance=0.8cm,
  box/.style={rectangle, draw, rounded corners=3pt,
    minimum width=11cm, minimum height=1.2cm,
    text width=10.5cm, align=center, font=\small},
  arrow/.style={-{Stealth[length=3mm]}, thick, gray!70}
]
\node[box, fill=prop1!10, draw=prop1!70] (p1)
  {\textbf{方向一：统一 Reward 框架}\\
   提供理论基础——self-consistency 与 verified reward 是同一连续谱};
\node[box, fill=prop2!10, draw=prop2!70, below=of p1] (p2)
  {\textbf{方向二：Meta-Capability 假说}\\
   定义训练目标——不训练具体知识，训练``如何运用知识''};
\node[box, fill=prop3!10, draw=prop3!70, below=of p2] (p3)
  {\textbf{方向三：对比式自我诊断}\\
   提供训练信号来源 A——模型自身的错误/正确轨迹对};
\node[box, fill=prop4!10, draw=prop4!70, below=of p3] (p4)
  {\textbf{方向四：多视角稳健性验证}\\
   提供训练信号来源 B——跨 framing 的一致性/不稳定性};

\draw[arrow] (p1) -- (p2);
\draw[arrow] (p2) -- (p3);
\draw[arrow] (p2) -- (p4);
\end{tikzpicture}
\end{center}

\noindent
如果资源有限需要优先排序：
\begin{itemize}[leftmargin=2em]
  \item \textbf{最高 risk/reward}：方向二（meta-capability 假说）——
    若验证通过则是范式转变，若失败则提供有价值的 negative result。
  \item \textbf{最可行的近期实验}：方向三（对比式自我诊断）——
    基础设施需求最低，可在单机上实现完整循环。
  \item \textbf{理论贡献最大}：方向一（统一 reward 框架）——
    提供覆盖全领域的统一理论。
  \item \textbf{最具原创性}：方向四（多视角稳健性）——
    唯一不依赖外部信号的训练范式，
    哲学上最深刻。
\end{itemize}

\end{document}
